% ---
% RESUMOS
% ---

% RESUMO em português
\setlength{\absparsep}{18pt} % ajusta o espaçamento dos parágrafos do resumo
\begin{resumo}

Este trabalho propõe uma análise comparativa entre ferramentas tradicionais de Static Application Security Testing (SAST) e ferramentas empresariais de desenvolvimento assistido por Inteligência Artificial para detecção de vulnerabilidades em aplicações Python. A pesquisa considera três abordagens: SAST tradicional puro, IA pura e uma abordagem híbrida SAST + IA. Serão avaliadas ferramentas como Bandit, Semgrep Community Edition, Cursor e OpenAI Codex, utilizando benchmarks e repositórios vulneráveis em Python. A metodologia envolve revisão bibliográfica, definição de protocolo experimental, execução das ferramentas e comparação dos resultados obtidos. A avaliação será feita por métricas quantitativas, como precisão, recall, F1-score e falsos positivos, além de critérios qualitativos, como clareza da explicação, severidade e utilidade da recomendação. Espera-se identificar em quais cenários a IA pode substituir, complementar ou enriquecer o uso de ferramentas SAST tradicionais em processos de segurança de software.

 \textbf{Palavras-chaves}: Static Application Security Testing; SAST; Inteligência Artificial; segurança de aplicações; vulnerabilidades em Python; agentes de codificação; DevSecOps; Bandit; Semgrep; OpenAI Codex; Cursor.
\end{resumo}

% ABSTRACT in english
\begin{resumo}[Abstract]
 \begin{otherlanguage*}{english}
This work proposes a comparative analysis between traditional Static Application Security Testing (SAST) tools and enterprise AI-assisted development tools for vulnerability detection in Python applications. The research considers three approaches: traditional SAST, pure AI-based analysis, and a hybrid SAST + AI approach. Tools such as Bandit, Semgrep Community Edition, Cursor, and OpenAI Codex will be evaluated using vulnerable Python benchmarks or repositories. The methodology includes a literature review, the definition of a standardized experimental protocol, tool execution, and comparison of the obtained results. The evaluation will consider quantitative metrics, such as precision, recall, F1-score, and false positives, as well as qualitative criteria, including explanation clarity, severity classification, and usefulness of recommendations. The study aims to identify in which scenarios AI can replace, complement, or enrich traditional SAST tools in software security processes.
   \vspace{\onelineskip}
 
   \noindent 
   \textbf{Keywords}: Static Application Security Testing; SAST; Artificial Intelligence; application security; Python vulnerabilities; coding agents; DevSecOps; Bandit; Semgrep; OpenAI Codex; Cursor.
 \end{otherlanguage*}
\end{resumo}